% !TeX root = ../nc-del-pezzo.tex
\section{有限対角と分離的な粗モジュライ}
\label{sec:separated-coarse-moduli}

\subsection{対称化した中心方程式と全同値の有限性}

以下では定理\ref{thm:zero-neighborhood-open-full-moduli}の開集合を用いる。
その全同値群oidは，定理\ref{thm:full-source-smooth-zero-neighborhood}の性質を保つ。
滑らかな可換零切断 $Z$ の全体を含む $W$ 上で，
全同値は通常の次数付き同型を次数スカラーで割ったものであり，
$s,t:\mathcal E_W\to W$ は滑らかである。
ambientの重み付き生成元の自由加群を $L$ とし，$e_1=6,e_2=4$ と書く。
既存の有限個の開条件により，$m\le e_d$ に対して
\[
 \operatorname{Sym}:(\operatorname{Sym}_{\mathcal O_W}L)_m
                  \xrightarrow{\sim} A_m
\]
は同型である。

\begin{lemma}[零方向も含む全同値の重み付き線型化]
\label{lem:all-graded-isom-weighted-linear}
補題\ref{lem:nonzero-graded-isom-weighted-linear}の生成元に関する結論は，
$W\times W$ 上の普遍的な通常次数付き同型スキーム上で成立する。
従って零方向・冪零方向を含む任意の基底変更後にも成立する。
特に全同値の各矢は，有効な重み付き線型変換 $g\in G_d$ を一意に定める。
\end{lemma}
\begin{proof}
未商化の通常次数付き同型スキームを $\widetilde{\mathcal E}$ と書く。
補題\ref{lem:graded-isom-affine-quotient}と
系\ref{cor:zero-neighborhood-full-graded-recognition}によって
$\widetilde{\mathcal E}\to\mathcal E_W$ は $\mathbb G_m$-torsorである。
$W$ は被約であり，$\mathcal E_W\to W$ は滑らかなので，
$\mathcal E_W$ と $\widetilde{\mathcal E}$ は被約である。
さらに，滑らかな射は開写像であり，$W\setminus Z$ は $W$ の稠密開集合なので，
$\widetilde{\mathcal E}^{\times}=s^{-1}(W\setminus Z)$ は
$\widetilde{\mathcal E}$ の稠密開集合である。
零切断の同値不変性により，この開上のtargetも非零方向にある。

普遍的同型の生成元の像を，次数2・3の完全対称化の基底で展開する。
一般には，次数2では $z$ の像に $q_2(x)$ が，次数1では
$y$ の像に $q_2(x)$，$z$ の像に
$\operatorname{Sym}(\ell_1(x)y)+q_3(x)$ が現れる。
これらの係数は $\widetilde{\mathcal E}$ 上の正則関数である。
非零方向の補題により，これらは稠密開上で消える。
被約性から全体で消える。
同じ議論で，次数1の $\gamma^2=\beta^3$ も全体で成立する。
対角係数の可逆性は，次数2では純粋な $x$ の完全対称化部分束，
次数1の次数3ではさらに $\operatorname{Sym}(xy)$ を加えた部分束を使い，
同型がそれらを可逆に保つことと，対応する階数1の商から従う。
ここで全ての低次数の積からなる部分加群を取るのではない。
次数スカラーで割ると，得た重み付き線型変換は一意な $G_d$ の元となる。
この全ては普遍スキーム上の恒等式であるから，非被約な基底変更にも保たれる。
\end{proof}


\begin{lemma}[対称化した中心元の可換模型]
\label{lem:commutative-central-shadow}
全可換零切断を保って $W$ をさらに不変開へ縮めると，
\[
 F^{\mathrm{sh}}=\operatorname{Sym}^{-1}(\Omega),\qquad
 C^{\mathrm{sh}}=\operatorname{Sym}_{\mathcal O_W}L/(F^{\mathrm{sh}})
\]
によって滑らかな可換del Pezzo族
\[
 X^{\mathrm{sh}}=\operatorname{Proj}_W C^{\mathrm{sh}}\longrightarrow W
\]
が定まる。零切断上では元の可換族に一致する。
全 $\mathcal O$ 保存同値の矢は自然にこの可換族の同型を誘導し，
従って群oidと両立する射
\[
 \iota:\mathcal E_W\longrightarrow
 J:=\underline{\operatorname{Isom}}_{W\times W}
        (X^{\mathrm{sh}}_1,X^{\mathrm{sh}}_2)
\]
がある。
\end{lemma}
\begin{proof}
$F^{\mathrm{sh}}$ は正則な係数を持つ重み $e_d$ の可換多項式で，
$z^2$ 係数は $1/2$ である。次数1では $y^3$ の係数が零切断上で
$1/3$ なので，その可逆な開集合へ制限する。
このとき重み付き射影空間の特異な座標点を避ける。
$z$ に関するmonic二次式による除算から $C^{\mathrm{sh}}$ は
残る変数の多項式環上で基底 $1,z$ を持ち，各重み成分は基底上有限局所自由である。
従って相対Projも基底上平坦で，有限表示かつ射影的である。
可換超曲面の非smooth locusの射影的な像を除けば，
零切断を失わず，$X^{\mathrm{sh}}\to W$ を滑らかにできる。
この開条件は重み付き線型変換に不変である。
重み付き超曲面の随伴公式とHilbert多項式から，
$-K=(n-1)H$，$H^n=d$ が成り立つ。

次数2では $z$ の平方完成を行い，次数1ではこれに加えて
$y$ の三次式の二次項を除く。
最高係数は可逆なので，fppf局所的な単元の根を使って通常の正規形に直せる。
従ってこの族は，入力文献\cite[定理8.8]{PoissonDP}の
内在的な可換モジュライ $\mathcal D_{n,d}$ の対象である。

前補題により，任意の次数付き同型は重み付き線型変換による。
ambientの関係式の次数は $e_d$ 未満なので，$B_m=A_m$ ($m<e_d$) から，
この変換はambientの同型にも持ち上がる。
中心関係式のイデアルの次数 $e_d$ 成分を比較すると
$\Omega_1$ は $\Omega_2$ の単元倍へ移る。
完全対称化は重み付き線型変換に可換するので，
$F^{\mathrm{sh}}_1$ は $F^{\mathrm{sh}}_2$ の同じ単元倍へ移る。
これが可換次数付き環の同型，従って $\operatorname{Proj}$ の同型を与える。
次数スカラーは $\operatorname{Proj}$ 上恒等であり，
構成は合成・逆射・基底変更と両立する。
\end{proof}

\begin{theorem}[全可換零切断を含む近傍の有限対角]
\label{thm:zero-neighborhood-finite-diagonal}
前補題の $W$ において，$\iota:\mathcal E_W\to J$ は閉埋込みである。
従って
\[
 (s,t):\mathcal E_W\longrightarrow W\times W
\]
は有限であり，$\mathfrak U_{n,d}=[W/\mathcal E_W]$ は
有限対角を持つ分離的な有限型DMスタックである。
\end{theorem}
\begin{proof}
\emph{古典同型層の有限性。}
可換del Pezzoモジュライ $\mathcal D_{n,d}$ の対角は有限である
\cite[定理9.2，\texttt{thm:affine-coarse-moduli}]{PoissonDP}。
従って，その対角を上の二つの可換族によって基底変更した $J$ は
$W\times W$ 上有限である。
この入力では，滑らかな正規形係数領域は判別式の主開でありaffine，
有効座標変換群は還元的，その全stabilizerは有限である。
その作用射の有限性はLunaの結果\cite[I, \S4, Corollaire, p.~90]{Luna1973}
による。非可換族の全同値についてこの結果を直接適用してはいない。

\emph{有限個の閉条件。}
古典同型は基本偏極のクラスを保つ。
そのためfppf局所的には，切断環 $C^{\mathrm{sh}}_1,C^{\mathrm{sh}}_2$ の
次数付き同型として一意的に次数スカラーを法として表せる。
これは可換正規形の全族の認識\cite[定理8.8]{PoissonDP}による。
従って
\[
 \widetilde J=\underline{\operatorname{Isom}}_{\mathrm{gr}}
          (C^{\mathrm{sh}}_1,C^{\mathrm{sh}}_2)\longrightarrow J
\]
は $\mathbb G_m$-torsorである。

$\widetilde J$ 上の普遍的な可換同型を生成元に展開し，
前補題で消えた $q_2,\ell_1,q_3$ の係数を零とする閉条件を課す。
次数1ではさらに $\gamma^2=\beta^3$ を課す。
残る変換は可逆な重み付き線型変換である。
この変換とその逆が，それぞれ $B_1,B_2$ の有限個の斉次定義関係式を零へ送る条件を課す。
各条件は有限局所自由な $B_m$ の中の有限個の係数の消滅であり，閉条件である。
得られた閉部分スキームは，その関手的定義から
$\underline{\operatorname{Isom}}_{\mathrm{gr}}(B_1,B_2)$ に一致する。
正逆の生成元写像が可逆であることは $\widetilde J$ の中で既に保証されている。
逆に非可換同型は前補題によってこの閉部分に入るので，条件の取りこぼしはない。
全ての条件は次数スカラーに不変であり，忠実平坦降下によって
$\mathcal E_W\to J$ は閉埋込みとなる。
有限射との合成は有限であり，商の対角を表示 $W\times W$ に引き戻すと
$(s,t)$ になるので，対角の有限性が降下する。
\end{proof}

\begin{proposition}[可換模型への表現可能な射]
\label{prop:representable-commutative-shadow}
定理\ref{thm:zero-neighborhood-finite-diagonal}の近傍上で，
対称化した中心方程式は自然な射
\[
 \operatorname{sh}:\mathfrak U_{n,d}\longrightarrow\mathcal D_{n,d}
\]
を定める。この射は代数空間によって表現可能であり，
零切断上では元の可換模型を返す。
\end{proposition}
\begin{proof}
補題\ref{lem:commutative-central-shadow}は，表示 $W$ 上に
$\mathcal D_{n,d}$ の族を与え，全同値群oidの矢に
合成・逆射・基底変更と両立する可換模型の同型を割り当てる。
可換族の降下によって，この構成は商 $\mathfrak U_{n,d}$ からの射となる。
定理\ref{thm:zero-neighborhood-finite-diagonal}の閉埋込み
$\mathcal E_W\hookrightarrow J$ は任意の基底上で単射である。
従って $\operatorname{sh}$ の相対自己同型の核は自明であり，
代数スタックの表現可能性判定
\cite[Lemma 101.6.2, Tag 04YY]{StacksProject}から結論を得る。
零切断上の一致は $F^{\mathrm{sh}}=F$ による。
\end{proof}

\begin{remark}
射 $\operatorname{sh}$ は，ここで構成した中心元と完全対称化に付随して，
非可換対象から滑らかな\emph{可換多様体}を取り出す。
第\ref{sec:poisson-shadow-comparison}節では，最低次数の交換子と組み合わせて
この可換模型上のPoisson構造を復元し，Poissonモジュライへの
stabilizer preservingな\'{e}tale射へ強める。
大域的な開埋込みには同節で定める一般次数の決定が残る。
\end{remark}

\begin{corollary}[零切断を含む分離的な粗モジュライ]
\label{cor:zero-neighborhood-separated-coarse}
定理\ref{thm:zero-neighborhood-finite-diagonal}の $\mathfrak U_{n,d}$ は，
分離的で整な有限型代数空間 $M^{\mathrm{nc}}_{n,d}$ を粗モジュライとして持つ。
次元は $\mathfrak U_{n,d}$ と同じである。
\end{corollary}
\begin{proof}
有限対角から慣性は有限である。
Keel--Mori定理\cite[Tag 0DUT]{StacksProject}により粗空間が存在し，
有限型性と分離性は\cite[Section 106.14, Tag 0DUW]{StacksProject}による。
既約性は粗射の全射性から，被約性は局所商の不変環の記述から従う。
有限慣性の粗射は幾何学的同型類を一点に移すため，次元も保つ。
\end{proof}

\begin{corollary}[安定化群を保つsliceと一般安定化群]
\label{cor:stabilizer-preserving-slice-generic-trivial}
定理\ref{thm:zero-neighborhood-finite-diagonal}の開近傍を用いると，
系\ref{cor:finite-etale-slice-near-classical}の
$[T/\Gamma]\to\mathfrak U_{n,d}$ は，
可換点を保った $\Gamma$ 不変affine縮小により
全点でstabilizer preservingにできる。
さらに $\mathfrak U_{n,d}$ の一般安定化群は自明であり，
非零方向には稠密な滑らかな代数空間の開部分が存在する。
\end{corollary}
\begin{proof}
有限慣性はproperである。
representable unramified射のstabilizer preserving locusは，
targetの慣性がproperなら開である\cite[Tag 0DUA]{StacksProject}。
可換点で既にstabilizer同型なので，この開集合を取り，
系\ref{cor:finite-etale-slice-near-classical}と同じ有限群のnormによって
$\Gamma$ 不変affine開へ縮める。

一般の可換模型 $X$ には $\Gamma=\mu_2$ しかなく，
この群は古典係数空間には自明，Poisson方向には符号で作用する
\cite[命題9.4と注意8.5]{PoissonDP}。
従って $T\subset(f+N)\times C_{n,d}$ の非零方向の部分では
$\Gamma$ のstabilizerは自明である。
この部分は $T$ の非空開集合であり，その像も $\mathfrak U_{n,d}$ の非空開集合である。
stabilizer preserving性から，像の全安定化群も自明である。
自明慣性を持つDMスタックは代数空間であり，非零方向では滑らかである。
$\mathfrak U_{n,d}$ は整なので，この開は稠密である。
\end{proof}

\begin{corollary}[粗空間の代数的局所模型]
\label{cor:algebraic-coarse-local-model}
前系のstabilizer preservingなsliceを用いると，粗空間の射
\[
 T/\Gamma_X\longrightarrow M^{\mathrm{nc}}_{n,d}
\]
は\'{e}taleであり，粗射との正方形
\[
\begin{array}{ccc}
[T/\Gamma_X]&\longrightarrow&\mathfrak U_{n,d}\\
\downarrow&&\downarrow\\
T/\Gamma_X&\longrightarrow&M^{\mathrm{nc}}_{n,d}
\end{array}
\]
はCartesianである。
特に，可換点の粗像を $[X]_{\mathrm c}$ と書けば，
\[
 \widehat{\mathcal O}_{M^{\mathrm{nc}}_{n,d},[X]_{\mathrm c}}
 \simeq R_X^{\Gamma_X}
\]
となる。この点で粗空間は正規，Cohen--Macaulayかつ解析的に一分枝である。

$\Gamma_X=\mu_2$ の場合には，この\'{e}tale局所模型は
\[
 \mathbb A^{b_{n,d}}\times
 \operatorname{Spec}\bigl(\mathbb C[C_{n,d}]^{(2)}\bigr)
\]
の原点を含む開部分である。
従って系\ref{cor:generic-nc-formal-sign-quotient}の接空間次元と特異性は，
実際の粗空間の代数的局所環についての結論でもある。
\end{corollary}
\begin{proof}
有限群商 $[T/\Gamma_X]$ は分離的であり，標的の対角も有限なので，
sliceの射はseparatedである。
有限慣性を持つスタックへのseparated，\'{e}tale，
stabilizer preservingな射は粗空間上の\'{e}tale射からの基底変更である。
これはwell-nigh affineな標的に対する
\cite[Lemma 106.13.7, Tag 0DUS]{StacksProject}を，
\cite[Lemma 106.13.8, Tag 0DUE]{StacksProject}の
stabilizer preservingな\'{e}tale被覆上で適用し，
Keel--Moriの構成における粗空間の\'{e}tale降下を行ったものである。
$T$ はaffineなので，その粗商は通常の有限不変式商 $T/\Gamma_X$ である。
これで最初の主張を得る。

sliceの中心点 $x$ は $\Gamma_X$ の固定点であり，
$\widehat{\mathcal O}_{T,x}\simeq R_X$ は同変な同型である。
有限群の平均化により不変式と完備化は可換するので，
$T/\Gamma_X$ の対応する点での完備局所環は $R_X^{\Gamma_X}$ である。
上の\'{e}tale射は両点で剰余体 $\mathbb C$ を保つので，
粗空間の完備局所環もこれに等しい。
系\ref{cor:formal-coarse-normal}の正規Cohen--Macaulay整域性を適用し，
excellentな局所環とその完備化の関係から実際の局所環の主張を得る。

$\Gamma_X=\mu_2$ のとき，係数方向 $N_X$ への作用は恒等で，
錐方向への作用は符号である。
$T$ は $(f+N_X)\times C_{n,d}$ の不変開集合であり，
有限商は不変閉集合を閉集合へ送るので，
$T/\Gamma_X$ は表示した積の開部分となる。
一般の $\Gamma_X$ の場合と異なり，この場合には係数方向との積も保たれる。
\end{proof}
